\documentclass[a4paper]{article}
\usepackage{amsfonts}
\usepackage{amsmath}
\usepackage{amssymb}
\usepackage{graphicx}     

\begin{document}
  
\noindent \large Auteur: Anna Voronovska \\
\noindent \large Studierichting: Informatica\\
\noindent \large Oefeningenreeks 6F nr. 16\\

\medskip

\normalsize

Puntsgewijze limiet:\\

$\lim\limits_{n \to +\infty} f_n(x) = \lim\limits_{n \to +\infty} \left(\dfrac{n^2+1}{n^2}\right)x = \lim\limits_{n \to +\infty} \left(1 + \dfrac{1}{n^2}\right)x = 1 \cdot x = x$\\

Dus de puntsgewijze limiet is $f(x) = x$ op $[0,1]$.\\

Uniforme limiet:\\

$\displaystyle\sup_{x \in [0,1]} \left|f_n(x) - f(x)\right| = \displaystyle\sup_{x \in [0,1]} \left|\left(\dfrac{n^2+1}{n^2}\right)x - x\right| = \displaystyle\sup_{x \in [0,1]} \dfrac{1}{n^2} x $\\

$= \dfrac{1}{n^2} \displaystyle\sup_{x \in [0,1]} x = \dfrac{1}{n^2}$\\

Dus:\\

$\lim\limits_{n \to +\infty} \displaystyle\sup_{x \in [0,1]} \left|f_n(x) - f(x)\right| = \lim\limits_{n \to +\infty} \dfrac{1}{n^2} = 0$\\

Omdat de supremum naar $0$ gaat, convergeert $f_n$ uniform naar $f(x) = x$ op $[0,1]$.



\begin{thebibliography}{999}
%\bibitem{a} Referentie
\end{thebibliography}
\end{document} 
